\subsection{An angular induction with curved narrow sets}

Besides a neighborhood of an affine codimension-two plane, we must
exclude a small quadratic sublevel set inside a hyperplane.  The latter
condition will detect the second fundamental form of a polynomial wall.
All frequency bins in the following lemma are bins of the original
unrotated grid.  Selecting a bin means selecting its entire field.

\begin{lemma}[Angular induction]\label{w:angular}
Assume the outer induction hypothesis in \Cref{w:weak} and the active
range \eqref{w:eq-active}.  Fix a subarray of energy $E_0\leq1$.
In every cube of a fixed spacetime grid of side $R^4$, prescribe some
of its frequency bins of side $1/R$.  Suppose their centers are covered
by $O((1+\log X)^{C_0})$ sets, each of one of the following types:
\begin{enumerate}[label=\textup{(\alph*)}]
\item a strip of width $O(R^{-1})$ about an affine $(n-2)$-plane;
\item the $O(R^{-1})$-neighborhood of
\[
 \{v\in H\cap B: |Q_H(v)|\leq R^{-4}\},
\]
where $H$ is an affine hyperplane, $B$ is a fixed bounded region,
and $Q_H$ has coefficient norm one in orthonormal coordinates on $H$
with bounded origin.
\end{enumerate}
The cover and the prescribed bins may depend on the macro cube.  Let
$U_{\mathrm{sel}}$ denote the resulting prescribed-bin sum on that cube.
For any fixed $c>0$,
\begin{equation}\label{w:eq-angular}
 \mu\{|U_{\mathrm{sel}}|\geq c\lambda\}
 \leq C C_{\mathrm{ind}}R^{-c_2}
       mL^{2n}g^{-p}E_0^{p/2}+C_NmL^{-N}.
\end{equation}
Here $N$ can be any prescribed fixed exponent, given sufficiently many
initial profile derivatives.  The exponent $c_2>0$ depends on the
dimensional choices in \eqref{w:eq-ab}, not on $u$ or $C_0$.
\end{lemma}

\begin{proof}
We give the geometric iteration, the measure transport, and the
exponent calculation separately.

\paragraph{The narrow geometry.}
Assign every selected bin to one covering set.  In a set of type
\textup{(b)}, choose one approximating point for each assigned bin and
sort the sizes of $\abs{\nabla Q_H}$ there into dyadic classes $m_1$,
down to $1/R$; put all smaller gradients in the last class.  There are
$O(\log R)$ classes.  Each assignment concerns the whole bin, not a
subset of its packet labels.

Consider one class through the $J$ iterations of ratio $K$.  If $m_1$
is a sufficiently small constant, the normalization of $Q_H$, its
bounded coefficients, and its small value at an approximating point
force its quadratic part to have norm bounded below.  Some component
of $\nabla Q_H$ then has a linear part bounded below.  The condition
$\abs{\nabla Q_H}\lesssim m_1$ puts this class in a strip of width
$O(m_1+R^{-1})$ within the $O(R^{-1})$-neighborhood of $H$.
For parent inverse widths
\[
 s\leq \frac{c}{Km_1},
\]
this is a codimension-two strip at the next child precision.

For a dyadic class with a positive lower gradient bound, the other
regime is
\[
 \frac{CK}{m_1}\leq s\leq \frac{R}{K}.
\]
Choose an approximating point in the parent.  Taylor expansion around
it places all assigned descendants within
\begin{equation}\label{w:eq-quadraticwidth}
 C\bigl(m_1^{-1}(s^{-2}+R^{-4})+R^{-1}\bigr)
 \lesssim (Ks)^{-1}
\end{equation}
of the tangent affine $(n-2)$-plane to its level set in $H$.
The last, small-gradient class never needs this second regime.  If the
gradient is bounded below by a fixed constant, only a bounded number
of initial steps are outside the second regime.  In general the gap
between the two regimes is a factor $O(K^2)$, so at most a bounded
number of ratio-$K$ steps need a crude estimate.

Use the functional $\mathcal T_b$ of \Cref{d:transport}, with
$P=2/(1-b)$.  Since $b<2/n$, we have $P\leq q_{n-2}$.
At every good step \Cref{d:strip-decoupling} applies after rescaling the width in
\eqref{w:eq-quadraticwidth} to $O(K^{-1})$.  The exceptional steps use
Cauchy--Schwarz.  Consequently the whole angular iteration, excluding
the transport factors, costs
\begin{equation}\label{w:eq-Kcost}
 C_\epsilon K^C R^\epsilon,
\end{equation}
where the exponent $C$ is independent of $J$.  Constants depending on
the fixed $J$ are allowed.  This distinction is essential: good-step
losses multiply to $R^\epsilon$, rather than $K^{CJ}$.

For this nonrecursive calculation one may smoothly truncate compact
profiles in native spacetime Fourier variables at radius $L^{1/4}$.
With sufficiently many fixed derivatives, every relevant subfield and
rescaling has error smaller than any prescribed negative power of $L$.
There are only polynomially many labels of packets meeting the domain.
The choice of $\kappa$ makes the fixed hierarchy $R=K^J$ small enough
relative to $L$ that the truncated supports satisfy the thickness and
overlap hypotheses of the packet estimates at each step.  Supremum
estimates for these approximations use ordinary convolution.  At the
terminal step we return to the original compact child fields, so none
of this truncation changes the induction's profile class.  The absolute
error estimate is justified below, including when the tested mass is
small.

\paragraph{Initial measure functional and recombination.}
Let $\nu$ be the original measure restricted to the set on the left
of \eqref{w:eq-angular}, and write $M_0=\nu(\R^{n+1})$.
For each cover/class index, pad the enumeration where needed and
restrict $\nu$ separately to the macro cubes.  On a unit cube $q$,
the density bound gives $\mu(q)^{1-b}\lesssim(m\chi)^{1-b}$.
H\"older on the constituent unit cubes, followed by the weighted
supremum estimate for the bandlimited approximation, bounds
$\lambda^2M_0$ by the sum of the initial $\mathcal T_b$ functionals,
multiplied by
\[
 C(m\chi)^{1-b}(1+\log X)^C.
\]
The sum of unit-cube convolution weights inside a $K^2$-cube is bounded
by the enlarged Schwartz weight used for that cube.  Thus the weighted
versions of the transport inequality apply without a factor counting
unit cubes.  During each calculation the angular subset prescribed on
the original macro is kept fixed, including in every shifted norm.

After $J$ steps, the transport factors and \eqref{w:eq-Kcost} give
\begin{equation}\label{w:eq-angular-factor}
 C_\epsilon K^C R^{\epsilon+2-(n+2)b}.
\end{equation}
We explain why the macro restrictions can now be recombined.  Write
\[
 A_{K,v}(x,t)=((x-2vt)/K,t/K^2).
\]
For a fixed angular chain, class index, and complete grid-shift
sequence, every step is a translated pushforward under $A_{K,v_j}$.
The grids on the right of the transport inequalities are global and
unrotated; rotations used to prove a strip inequality do not change
these lists.  Hence the composition is
\begin{equation}\label{w:eq-affine-composition}
 \tau_{z_{\rm sh}}\circ A_{R,V},\qquad
 V=v_1+v_2/K+\cdots+v_J/K^{J-1},\qquad |V|\leq C_n,
\end{equation}
with a translation $z_{\rm sh}$ independent of the original macro cube.
The inverse image of a terminal $K^2$-cube has diameter at most
$C_nR^2K^2\leq C_nR^4$, regardless of the size of $z_{\rm sh}$.  It therefore
meets only boundedly many original macros.  For their transported
measures this gives, on every terminal cube $Q$,
\[
 \sum_{\mathcal M}\nu_{\mathcal M}(Q)^b
 \leq C_n\left(\sum_{\mathcal M}\nu_{\mathcal M}(Q)\right)^b.
\]
The combined measure is dominated by one translated pushforward of
$\nu$, has mass at most $M_0$, and has ball parameter
\begin{equation}\label{w:eq-childm}
 m'\asymp C mR^{n+1}.
\end{equation}
Infinite shift sequences are summed with their summable weights;
neither the mass, the ball constant, nor the number of occupied cubes
depends on their absolute translation.  Since assignments were by
whole original $1/R$-bins, the terminal field is the same full compact
bin field, or zero, on every contributing macro.  Thus the outer
induction will charge each original bin energy once.

\paragraph{The terminal induction.}
Put $l=L/R$.  In each terminal $K^2$-cube with nonzero original child
supremum, move its mass to a point arbitrarily close to that supremum
and bound its $P$-norm by the supremum.  This changes ball bounds by
at most $K^C$.  Sort the cube masses as $\nu_{\mathrm{child}}(Q)\asymp
m'd$.  Unit covering of the inverse image, whose volume scales as
$R^{n+2}$ while \eqref{w:eq-childm} scales as $R^{n+1}$, gives
\begin{equation}\label{w:eq-dcap}
 d\leq CK^C R\chi.
\end{equation}
The moved measure has ball parameter $CK^Cm'$, relative unit density
at most $\min(1,Cd)$, and, after fixed constant inflation, admissible
moment parameter $D=1$.  Its relevant time range is that of the
original compact child field.  Spatial boxes can be partitioned as
in the theorem.

Apply the outer induction at $l$ and integrate the weak $L^p$ bound
against the trivial mass bound $M_0$.  Since $p>2$, the resulting
$L^2$ integral is bounded by
\begin{equation}\label{w:eq-terminalL2}
 CK^C C_{\mathrm{ind}}^{2/p}(m')^{2/p}
 l^{4n/p-2e}d^{2a}M_0^{1-2/p}E_{\mathrm{bin}},
\end{equation}
where $E_{\mathrm{bin}}$ is that bin's original coefficient energy.
Indeed, integrating a distribution estimate
$\nu\{|V|>t\}\leq\min(M_0,At^{-p})$ gives
$\int|V|^2\,\dd\nu\leq C_p A^{2/p}M_0^{1-2/p}$.
Multiplication by $(m'd)^{b-1}$ returns to the mass-power functional;
the fixed $K$-volume factor costs another $K^C$.  The dyadic $d$ sum
converges at its lower end because
\[
 b-1+2a=\delta>0.
\]
By \eqref{w:eq-dcap}, its total cost is
$C(K^CR\chi)^\delta$.

To justify returning from the Fourier approximations even on cubes
with zero original supremum, normalize $m=1$.  There are polynomially
many occupied cubes, and concavity gives
\[
 \sum_Q\nu(Q)^b\leq(\#Q)^{1-b}M_0^b.
\]
Choosing the fixed truncation order large enough makes the entire
error, with all prefactors restored, at most $L^{-N'}M_0^b$ for any
prescribed $N'$.  Either this absorbs into $\lambda^2M_0$, or, using
$\lambda\geq L^{-C}$ from \eqref{w:eq-active},
\[
 \lambda^2M_0\lesssim L^{-N'}M_0^b
 \quad\Longrightarrow\quad
 M_0\lesssim L^{-(N'-2C)/(1-b)}.
\]
The latter is the asserted absolute negligible mass after choosing
$N'$ large.  Restoring homogeneity supplies the factor $m$.

\paragraph{The gain.}
Sum the disjoint bin energies.  The factors of $m$ and $\chi$ before
solving for $M_0$ are exactly $m^{2/p}\chi^{2a}$.  Apart from
$K^CR^\epsilon$ and fixed logarithmic powers, the exponent of $R$ is
\begin{equation}\label{w:eq-Gamma}
 \Gamma(p,\delta)=2-(n+2)b+(n+1)(b-1+2/p)-4n/p+2e+\delta.
\end{equation}
At $p=2$, $\delta=0$, direct simplification gives
$\Gamma=\beta-b<0$.  Solving for $M_0$ multiplies this exponent by
$p/2>0$.  First choose $\delta$ and $p_{\max}-2$ small, then $J$
large so that $K^C=R^{C/J}$ is paid by the strict margin, and finally
choose the decoupling loss small.  The remaining logarithmic factors
are absorbed by increasing $B$ for the fixed $C_0$.  This proves
\eqref{w:eq-angular} with a fixed $c_2>0$.
\end{proof}
